%% Hand-written B blocks of fork-new items (no parent B block and no parent generated signature, binders or hashes).
%% Read by mkB.py: each chunk starts with a line `%% fork-new after <anchor ID>` and holds exactly one item block.
%% fork-new after I.adele.thm:adele-fundamental-domain

\begin{theorem}[$\Q\backslash\A$ is compact]
\tfstmt{I.adele.thm:rat-quotient-compact-top}
\lean{GL2TF.rat_quotient_compact_top}
\uses{I.adele.def:adele-ring,I.adele.def:zhat,I.adele.def:finite-adeles-model,I.adele.def:rat-embedding,I.adele.lem:primes-fact-instance}
\end{theorem}
\begin{tfformal}{I.adele.thm:rat-quotient-compact-top}
\tfhyps{I.adele.thm:rat-quotient-compact-top}{}
\begin{tffidelity}
Fork-new item (added in the algebraic\_jl fork, not in the parent at 84818d0): there is no generated signature, binder
row or sem\_hash for it, and it carries no audit. Lean statement: \texttt{CompactSpace (Adele ⧸ (algebraMap ℚ
Adele).range.toAddSubgroup) ∧ T2Space (Adele ⧸ (algebraMap ℚ Adele).range.toAddSubgroup)}. ``$\A/\Q$'' is the quotient
of the additive group \texttt{Adele} by the range of \texttt{algebraMap ℚ Adele}, with Mathlib's quotient topology. No
binders. It is the measure-free part of the parent's \texttt{I.adele.thm:rat-quotient-compact} (dropped in the fork,
since its other clauses normalize a Haar measure), extended by the Hausdorff clause. \texttt{\#print axioms}:
\texttt{propext}, \texttt{Classical.choice}, \texttt{Quot.sound}.
\end{tffidelity}
\end{tfformal}
\begin{proof}[Lean proof]
Compactness: the image of $[0,1]\times\hat\Z$ (compact by \texttt{isCompact\_Icc.prod isCompact\_zhat}) under the
quotient map is everything, because by \cref{I.adele.thm:adele-fundamental-domain} every $x$ is $q+y$ with $y\in
[0,1)\times\hat\Z$ and then $x$ and $y$ have the same class; a continuous image of a compact set is compact. Hausdorff:
the subgroup is closed (its carrier is \texttt{Set.range (algebraMap ℚ Adele)}, closed by the first clause of
\cref{I.adele.thm:rat-discrete}); Mathlib's instance \texttt{QuotientAddGroup.instT3Space} for the quotient by a closed
normal subgroup then gives \texttt{T2Space}.
\end{proof}
